\section{Literature Review}
\subsection{Introduction}

This chapter reviews short-term solar photovoltaic generation forecasting with machine learning, focusing on data, model choice and evaluation relevant to New South Wales (NSW). The retained literature did not identify a benchmark matching this project’s NSW1 market-cleared target and combination of reanalysis with completed-hour AEMO generation history.

The review draws on the project's July 2026 search notes. arXiv was queried through 15 Boolean expressions combining terms such as solar PV forecasting, machine learning, irradiance, LSTM, physics-guided and reanalysis. Semantic Scholar was queried through three broader expressions. Coverage ran from 2020 to 2026, with foundational pre-2020 papers retained where they established architectures still in use. Inclusion required machine learning or statistical forecasting of solar irradiance or PV generation with reported quantitative error metrics. Wind-only forecasting, papers without reported metrics, and hardware or policy analysis were excluded. Screening was conducted by a single reviewer, and Appendix C records the flow.

The bibliography holds 37 entries: 30 forecasting studies, six data or service sources and one method reference. Per-record screening decisions were not retained, so the review cannot be reproduced record by record. Its gap statement is therefore bounded by the retained material. Appendix C records the reconciliation.

\subsection{Data Characteristics and Preprocessing}

Two distinctions structure the data literature. Pfenninger and Staffell \cite{pfenninger2016longterm} evaluate hourly PV simulations based on reanalysis and satellite data in Europe. Ramirez Camargo and Schmidt \cite{ramirez2020simulation} compare ERA5-land with MERRA-2-based renewables.ninja output at 57 Chilean plants, finding only modest improvements under their modelling assumptions. These studies motivate historical data use without validating this NSW meteorological aggregation. By contrast, Nielsen et al. \cite{nielsen2021irradiancenet} use satellite sequences for short-term irradiance prediction. The second distinction is temporal: hourly meteorology \cite{renewables2025api} and five-minute AEMO dispatch \cite{aemo2024nemosis} require explicit alignment. AEMO's data model defines the two series as regional aggregated unconstrained intermittent generation forecast and regional aggregated semi-scheduled cleared MW \cite{aemo2026datamodel}. The cleared series is the dispatch outcome, so it sits at or below the unconstrained series, and the gap is a market-clearing effect rather than a weather effect. This project retains completed-hour dispersion summaries alongside hourly history.

Multi-source studies address different information needs. Rabault et al. \cite{rabault2025datafusion} combine reanalysis and satellite products to improve GHI estimates against Norwegian observations; Perera et al. \cite{perera2024dayahead} combine generation and weather across 101 Western Australian PV locations; Prasad and Kay \cite{prasad2021prediction} convert satellite-derived irradiance forecasts into power at Queensland and Victorian sites. Wen et al. \cite{wen2026multisource} combine historical, meteorological and forecast variables with cross-dimensional feature extraction. These complementary inputs do not test the AEMO history contrast used here, and their differing tasks cannot establish one universally sufficient spatial resolution.

\subsection{Algorithms and Modelling Approaches}

The literature spans statistical, boosting, recurrent, transformer and constrained methods, but added complexity must be judged within comparable forecasting tasks.

Architecture rankings depend on the task and comparison. Mukhoty et al. \cite{mukhoty2019seq2seq} report strong LSTM encoder-decoder performance against their forecasting alternatives, while Mollasalehi and Farhadi \cite{mollasalehi2025review} review model prevalence, complexity and execution time rather than establishing a universal accuracy-to-speed winner. Feng et al. \cite{feng2018clustering} report benefits from clustering and model blending, and Perera et al. \cite{perera2022multires} report a 3.81 percent average MASE reduction from forecast combinations against their best individual model. Abdullah and Mohammed \cite{abdullah2026physics} report GHI RMSE of 19.53 W/m2 versus 30.64 W/m2 for attention-based alternatives in Sudan. Their physical guidance includes engineered inputs, so this result does not validate the present non-negativity loss. Hoang \cite{hoang2025efficient} reports transformer R2 of 0.9696, whereas Scott et al. \cite{scott2025transformers} find near-baseline performance with poorer sunny-day results. SolarTformer \cite{basu2026solartransformer} and CT-PatchTST \cite{lu2025ctpatchtst} investigate attention and input relationships, while Sudharshan et al. \cite{sudharshan2026spatiotemporal} combine SARIMA and BiLSTM for seasonal irradiance prediction. The differing targets and protocols prevent a common model ranking.

Australian evidence is narrower. Boland et al. review wind and solar farm output forecasting for the National Electricity Market \cite{boland2022forecasting}, and Farah and Boland forecast individual Australian solar farms at five-minute resolution against clear-sky-index persistence and the model AEMO then used \cite{farah2021time}. Both address farm or portfolio output rather than a regional market-cleared aggregate.

\subsection{System and Architecture Considerations}

Operational constraints receive less attention than accuracy in the reviewed work. Dao et al. \cite{dao2020ensemble} consider ensembles on a Raspberry Pi, while Schubnel et al. \cite{schubnel2025solarcross} use a more demanding multimodal model. Their settings illustrate why deployment cost and available inputs belong alongside accuracy in model selection.

The staged baseline, tree and recurrent comparison tests incremental gains within the project's resources. Transformers are excluded by scope, not evidence of inferiority. The complexity claim in Abdullah and Mohammed \cite{abdullah2026physics} motivates comparison, but one study of engineered physical inputs is not evidence that the present loss penalty causes an improvement.

\subsection{Evaluation Metrics and Experimental Design}

Point-error metrics such as RMSE support model comparison, but do not describe the full distribution of future generation. Voyant et al. \cite{voyant2025nice} motivate evaluating more than one error characteristic. Probabilistic work is already present in the retained literature: SolarCrossFormer \cite{schubnel2025solarcross} produces probabilistic irradiance forecasts. The present study evaluates point forecasts and uncertainty in aggregate performance; its bootstrap intervals are not predictive intervals for individual future hours.

Evaluation choices limit comparisons across studies. Perera et al. \cite{perera2024dayahead} report a statistically significant improvement, so significance testing cannot be presented as absent from the reviewed literature. Pedro and Coimbra \cite{pedro2012assessment} explicitly include persistence and ARIMA, supporting the use of strong target-only baselines here. Scott et al. \cite{scott2025transformers} identify sunny-day degradation, while Prasad and Kay \cite{prasad2021prediction} show that delivery latency affects live forecasts. These findings motivate baseline comparison, paired uncertainty assessment and explicit input timing in the present benchmark. Irradiance and generation also remain distinct targets: Prasad and Kay \cite{prasad2021prediction} model their conversion, whereas this project directly forecasts market-cleared MW. Cross-study errors cannot therefore be treated as comparable effect sizes.

Verification practice is itself codified. Yang et al. set out the requirements for deterministic solar forecast verification \cite{yang2020verification}, and the ROPES guideline makes operational information availability a condition of a credible forecast claim \cite{yang2019guideline}. Scoring against persistence, or its clear-sky-index form, is standard \cite{farah2021time}, and the canonical chain from issue-time numerical weather prediction to PV power across the 1 to 24-hour band is described by Lorenz et al. \cite{lorenz2009irradiance}. Two cautions follow for a reanalysis-based design. Forecast models are often built ex post from observed weather over the target period while practice relies on numerical weather prediction \cite{yang2022historical}, and operational forecasts differ measurably from reanalysis over the same region and metric \cite{huang2018assessing}.

\subsection{Conceptual Framework}

The framework separates three nested information layers from model configuration. Weather and calendar variables form the base; completed-hour history adds the recent level and trajectory; and five-minute dispersion adds variation hidden by the hourly mean. Because history and dispersion describe the same completed hour, dispersion is expected to add more over meteorology alone, less once history is included, and less as the horizon lengthens.

Model family moderates extraction: trees and recurrent networks may use the same signal differently. Physical guidance can enter through engineered variables or the training loss \cite{abdullah2026physics}, but the constrained model changes architecture and loss together, so its ranking cannot isolate either mechanism. The primary estimand is the within-model error change when dispersion is added; architecture and ensemble comparisons are secondary. This follows multi-source practice \cite{perera2024dayahead}, \cite{prasad2021prediction}, \cite{rabault2025datafusion}, \cite{wen2026multisource} without assuming extra physical preprocessing improves accuracy \cite{voyant2025clearsky}.

\subsection{Research Gap and Synthesis}

The reviewed evidence supports testing data integration while treating architecture gains as task-dependent because targets, input availability and protocols prevent a universal ranking \cite{perera2024dayahead}, \cite{scott2025transformers}, \cite{mollasalehi2025review}. The retained studies do not match the present comparison of meteorology-only inputs, hourly generation history and AEMO dispersion features for regional market-cleared solar dispatch in NSW1. Perera et al. \cite{perera2024dayahead} instead forecast distributed rooftop generation in Western Australia. This source-and-target distinction defines an applied contribution, not priority over every use of AEMO data or a new method based on geography alone.

Searches of Crossref, OpenAlex and Semantic Scholar in September 2026 for DISPATCHREGIONSUM, SS\_\allowbreak{}SOLAR\_\allowbreak{}CLEAREDMW, SS\_\allowbreak{}SOLAR\_\allowbreak{}UIGF and semi-scheduled found no peer-reviewed forecast of this regional series. The closest Australian work forecasts individual farms or portfolios \cite{boland2022forecasting}, \cite{farah2021time}. AEMO documents the aggregate's forecasting and dispatch process \cite{aemo2022handbook} but publishes no matching benchmark.

Transfer is limited by different targets and input availability. Much reviewed work predicts irradiance rather than market-cleared MW. Retrospective reanalysis differs from weather forecasts available at operational issue times, so neither a fixed deployment penalty nor unchanged model rankings is established. The present meteorology is not independently validated against local ground-station observations; published validation in other regions \cite{pfenninger2016longterm}, \cite{ramirez2020simulation} motivates the source choice but does not establish local accuracy.

The study joins reanalysis and AEMO data for NSW, compares completed-hour history inputs, and adds paired significance tests. Distributed rooftop forecasting \cite{perera2022multires} remains outside this market-cleared target, which does not cover total NSW solar production.
